% Design Rationale — COUNTED in 15-page limit
% Justification of design decisions including STEEPLE considerations

\section{Design Rationale}\label{sec:rationale}

Every major design decision was evaluated against the seven STEEPLE dimensions (Social, Technological, Economic, Environmental, Political, Legal, Ethical). Technical trade studies then addressed hardware, software architecture, detection, path planning, and mission control.

\subsection{STEEPLE Analysis}\label{sec:steeple}

\begin{table}[H]
  \centering\compactTable
  \caption{STEEPLE assessment of the autonomous SAR drone system.}\label{tab:steeple}
  \begin{tabularx}{\linewidth}{@{}l X@{}}
    \toprule
    \textbf{Dimension} & \textbf{Key consideration and design response} \\
    \midrule
    Social & Serves community safety via volunteer rescue teams; operator-in-the-loop verification maintains public trust~\cite{murphy2014}. \\
    Technological & Pi~5 with TFLite (4 packages vs 90+); dual-backend architecture enables GPU development with edge deployment~\cite{edgeai, tflite}. \\
    Economic & \pounds565 total cost ($\sim$10\% of DJI Matrice 30T); onboard inference eliminates high-bandwidth comms and dedicated analyst. \\
    Environmental & Three-layer geofence + firmware fence protects adjacent SSSI; 30\,m waypoint buffer; \SI{3}{\kilo\gram} AUW minimises wildlife disturbance~\cite{wca1981}. \\
    Political & UK CAA Open Category A3 compliant~\cite{caa2024uas}; VLOS only (BVLOS requires SORA~\cite{jarus2019sora}); MIT licence (R12). \\
    Legal & Images stored locally (no cloud); SSSI compliance auditable via logged boundary distances at every loop iteration~\cite{wca1981}. \\
    Ethical & No autonomous landing without human confirmation; 120\,s auto-reject prioritises continued search over uncertain detections. \\
    \bottomrule
  \end{tabularx}
\end{table}

\noindent With the STEEPLE framework confirming project viability across all seven dimensions, the following subsections detail the technical trade studies that determined the system's operating parameters, hardware platform, software architecture, and mission control strategy.

\subsection{Search Parameter Optimisation}\label{sec:search-param-opt}

The search pattern is governed by three coupled decision variables---altitude~$h$, scan angle~$\theta$, and speed~$v$---that jointly determine five competing objectives: coverage completeness, detection probability, mission time, energy consumption, and NFZ safety margin.

These objectives cannot be optimised independently. Lowering altitude improves detection confidence but narrows the camera footprint and increases the number of scan lines. Increasing speed reduces mission time but shortens target dwell time and therefore detection probability.

\begin{figure}[H]
  \centering
  \includegraphics[width=0.78\linewidth]{coupling_matrix}
  \caption{Variable coupling matrix. The altitude--speed pair exhibits the dominant interaction, jointly determining coverage rate and detection probability. Off-diagonal cells identify which performance metric each variable pair affects.}
  \label{fig:coupling-body}
\end{figure}

A 216-configuration parametric sweep (6~altitudes $\times$ 36~scan angles, Appendix~\ref{sec:path-opt}) evaluated every combination on the real AENGM0074 polygon geometry using a physics-based energy model. The sweep reveals that scan angle selection dominates altitude selection for energy efficiency: varying the angle from optimal (\SI{70}{\degree}) to worst-case (\SI{0}{\degree}) increases energy by 100--200\%, whereas changing altitude from \SI{35}{\metre} to \SI{50}{\metre} saves only 30--40\%. A dependency chain links each parameter to a physical measurement or constraint:

\begin{enumerate}[nosep]
  \item \textbf{Target $\rightarrow$ altitude}: the mannequin's \SI{0.5}{\metre} width must subtend $\geq$10~model-input pixels for reliable detection, capping altitude at \SI{35}{\metre} with a $1.4\times$ margin.
  \item \textbf{Altitude $\rightarrow$ speed}: the \SI{24}{\metre} along-track footprint at \SI{35}{\metre}, combined with \SI{4.8}{FPS} inference, requires $\geq$10~frames per target, yielding $v \leq \SI{11.5}{\metre\per\second}$; the altitude-dependent schedule assigns $v = \SI{8}{\metre\per\second}$.
  \item \textbf{Polygon geometry $\rightarrow$ scan angle}: alignment with the longest edge (\SI{70}{\degree}) minimises U-turns from 9 to 5 and reduces energy from \SI{25.4}{Wh} to \SI{12.6}{Wh}.
  \item \textbf{Footprint width $\rightarrow$ NFZ margin}: the \SI{32}{\metre} footprint at \SI{35}{\metre} requires \SI{23}{\metre} Root Sum of Squares (RSS) standoff; the \SI{30}{\metre} buffer provides a $2.2\times$ safety factor.
\end{enumerate}

\begin{figure}[H]
  \centering
  \includegraphics[width=0.78\linewidth]{altitude_speed_tradeoff}
  \caption{Altitude--speed trade-off space. Colour encodes composite mission score (40\% detection $+$ 35\% coverage $+$ 25\% energy). The selected operating point (\SI{35}{\metre}, \SI{8}{\metre\per\second}) sits above the 95\% detection contour in the high-score region, trading \SI{4.4}{Wh} of energy (10.9\% of battery) for a comfortable detection margin.}
  \label{fig:alt-speed-body}
\end{figure}

\noindent The selected operating point (\SI{35}{\metre}, \SI{70}{\degree}, \SI{8}{\metre\per\second}) is not the energy minimum---that belongs to \SI{50}{\metre}/\SI{70}{\degree}/\SI{10}{\metre\per\second} at \SI{8.2}{Wh}---but it is the best Pareto compromise: 96\% coverage, $P_d > 99.97\%$, \SI{12.6}{Wh} (10.9\% of battery), and a \SI{50}{\metre} NFZ margin with $2.2\times$ surplus. The 4\% coverage gap lies within a fenced wildlife reserve with negligible casualty probability. The full derivation, energy model, strategy comparison, and sensitivity analysis are provided in Appendices~\ref{sec:path-opt} and~\ref{sec:design-strategy}.

\subsection{Parameter Derivation Chain}\label{sec:param-chain}

Every operating parameter traces back to a physical measurement or regulatory constraint through an unbroken eight-step chain (Table~\ref{tab:param-chain}). Full equations, RSS overlap analysis, and sensitivity analysis are provided in Appendices~\ref{sec:path-opt} and~\ref{app:decision-flow}.

\begin{table}[H]
  \centering\compactTable
  \caption{Parameter derivation chain: each value is derived from the previous step.}\label{tab:param-chain}
  \begin{tabularx}{\linewidth}{@{}c l c X@{}}
    \toprule
    \textbf{Step} & \textbf{Parameter} & \textbf{Value} & \textbf{Derivation} \\
    \midrule
    1 & Detection ceiling & \SI{63}{\metre} & Mannequin (\SI{0.5}{\metre} wide) must subtend $\geq$20 model-input pixels \\
    2 & Operating altitude & \SI{35}{\metre} & 30\% margin below ceiling (36\,px target, $1.8\times$ above minimum) \\
    3 & Max speed & \SI{11.5}{\metre\per\second} & \SI{24}{\metre} footprint $\div$ (10 frames / 4.8\,FPS); adopted \SI{8}{\metre\per\second} \\
    4 & Lane spacing & \SI{25.7}{\metre} & \SI{32.2}{\metre} footprint $\times$ 0.80; 20\% overlap gives $1.4\times$ margin over \SI{4.5}{\metre} RSS error~\cite{galceran2013survey, choset2001coverage} \\
    5 & Scan angle & \SI{70}{\degree} & Longest-edge alignment minimises U-turns (5 vs 9), saves \SI{13}{Wh} \\
    6 & NFZ buffer & \SI{30}{\metre} & Half-footprint (16.1) + GPS CEP (3.0) + stopping dist.\ (10.7) = 29.8\,m~\cite{wca1981} \\
    7 & Coverage & 96\% & 4\% gap in fenced wildlife reserve (negligible casualty probability) \\
    8 & Search energy & \SI{12.6}{Wh} & 10.9\% of \SI{115.4}{Wh} battery; 14\,min remaining for transit/verify/RTL \\
    \bottomrule
  \end{tabularx}
\end{table}

\subsection{Hardware Platform Selection}\label{sec:hw-rationale}

The hardware integrates a Hexsoon EDU-450 airframe, CubePilot Orange+ flight controller with ArduPilot firmware, Raspberry Pi~5 companion computer, and IMX296 global shutter camera.

\begin{itemize}
    \item \textbf{Global shutter camera.} The IMX296 was selected over cheaper rolling-shutter alternatives because rolling-shutter distortion during forward flight at up to \SI{10}{\metre\per\second} skews bounding boxes and degrades detection accuracy. The CSI interface leaves USB ports free for telemetry radios.
    \item \textbf{No thermal imaging.} Thermal cameras add \pounds2\,000--\pounds10\,000 to the platform cost and have lower spatial resolution than RGB at equivalent price points. This project deliberately explores RGB-only detection feasibility at the low-cost extreme~\cite{thermal_sar}.
    \item \textbf{MAVProxy bridge.} A Python~3.13 incompatibility with pyserial (dropped bytes causing \texttt{BAD\_DATA} errors) is resolved by interposing MAVProxy as a UDP bridge, which reads the serial port reliably using its own C-level handler~\cite{mavproxy}.
\end{itemize}

\subsubsection{Companion Computer Trade Study}

All four trade studies in this section use weighted multi-criteria decision analysis (MCDA). In each case, we established criteria weights before scoring candidates, based on mission requirements and operational constraints, to prevent post-hoc rationalisation.

Three candidate companion computers were evaluated (Table~\ref{tab:mcda-companion}). Cost and Python ecosystem compatibility are paramount for a student team, while raw GPU performance is less critical given the lightweight YOLOv8n model.

\begin{table}[H]
  \centering\compactTable
  \caption{MCDA: companion computer selection (1--5 scale, 5 = best)}
  \label{tab:mcda-companion}
  \begin{tabularx}{\textwidth}{X c c c c}
    \toprule
    \textbf{Criterion} & \textbf{Weight} & \textbf{Pi~5} & \textbf{Jetson Nano} & \textbf{Intel NCS2} \\
    \midrule
    Unit cost (\pounds)          & 0.20 & 5 (\pounds75)  & 3 (\pounds180) & 4 (\pounds100+host) \\
    Power consumption            & 0.10 & 4 (5--12\,W)   & 3 (5--20\,W)   & 3 (host+1.5\,W) \\
    GPIO / CSI interface         & 0.15 & 5 (40-pin+CSI) & 4 (40-pin+CSI) & 1 (USB only) \\
    Community \& documentation   & 0.20 & 5 (largest)    & 4 (strong)     & 2 (discontinued) \\
    Python~3.13 compatibility    & 0.20 & 5 (native)     & 3 (JetPack lag) & 2 (OpenVINO lag) \\
    TFLite / edge-AI support     & 0.15 & 5 (XNNPACK)   & 4 (TensorRT)   & 2 (OpenVINO only) \\
    \midrule
    \textbf{Weighted total}      &      & \textbf{4.90}  & 3.50           & 2.35 \\
    \bottomrule
  \end{tabularx}
\end{table}

\noindent The Pi~5 achieves the highest weighted total. Sensitivity analysis ($\pm$0.05 weight perturbation across all 12 scenarios) confirms the Pi~5 retains its rank in every case (minimum score 4.72 vs Jetson Nano maximum 3.68). The same perturbation was applied to all four MCDA tables; in each case the selected option retained its rank.

\subsubsection{Communication Architecture Trade Study}

Three architectures for companion-to-autopilot communication were compared (Table~\ref{tab:mcda-comms}).

\begin{table}[H]
  \centering\compactTable
  \caption{MCDA: communication architecture (1--5 scale, 5 = best)}
  \label{tab:mcda-comms}
  \begin{tabularx}{\textwidth}{X c c c c}
    \toprule
    \textbf{Criterion} & \textbf{Weight} & \textbf{Direct Serial} & \textbf{MAVProxy Bridge} & \textbf{ROS~2} \\
    \midrule
    Latency                    & 0.20 & 5 (direct)    & 4 (UDP hop)       & 3 (Data Distribution Service (DDS) overhead) \\
    Reliability                & 0.25 & 2 (py3.13 bug) & 5 (C-level read) & 4 (mature) \\
    Multi-consumer support     & 0.20 & 1 (exclusive) & 5 (N outputs)     & 5 (pub/sub) \\
    Python~3.13 compatibility  & 0.20 & 1 (broken)    & 5 (works)         & 3 (partial) \\
    Implementation complexity  & 0.15 & 5 (trivial)   & 4 (one process)   & 1 (full stack) \\
    \midrule
    \textbf{Weighted total}    &      & 2.65          & \textbf{4.65}     & 3.35 \\
    \bottomrule
  \end{tabularx}
\end{table}

\noindent Direct serial is the simplest option but is rendered non-viable by a Python~3.13 pyserial regression that drops bytes at 921\,600~baud. MAVProxy resolves this by reading the serial port with its own C-level handler and forwarding via UDP, simultaneously enabling multi-consumer access (scripts on port~14550, Mission Planner on TCP~5762). ROS~2 offers the richest middleware but introduces substantial deployment complexity disproportionate to the system's single-node architecture.

\subsection{Software Architecture Rationale}\label{sec:arch-rationale}

Four principles govern the architecture: (1)~\textbf{Modular separation}---each concern is isolated with a narrow interface (e.g.\ \texttt{detect\_in\_image(frame)} $\to$ \texttt{(found, x, y, conf)}); (2)~\textbf{Same code, simulation and real}---no conditional branches distinguish platforms; a \texttt{MODE} flag selects the camera, auto-detection selects the flight controller; (3)~\textbf{Configuration-driven tuning}---all parameters in \texttt{config.py}, transitioning to real hardware requires only parameter adjustment; (4)~\textbf{Dual inference backend}---auto-detects Ultralytics (laptop) or TFLite (Pi) at import, both producing identical $[1, 5, 8400]$ tensors.

\subsection{Detection Model Selection}\label{sec:model-rationale}

YOLOv8n (the nano variant, 3.2~million parameters, 8.7~GFLOPs) was selected after benchmarking the YOLO model family on the target hardware~\cite{yolov8}. Table~\ref{tab:mcda-model} summarises the weighted comparison.

\begin{table}[H]
  \centering\compactTable
  \caption{MCDA: detection model selection (1--5 scale, 5 = best)}
  \label{tab:mcda-model}
  \begin{tabularx}{\textwidth}{X c c c c c}
    \toprule
    \textbf{Criterion} & \textbf{Weight} & \textbf{YOLOv8n} & \textbf{YOLOv8s} & \textbf{YOLOv8m} & \textbf{SSD MobileNet} \\
    \midrule
    Inference speed (Pi~5 CPU)   & 0.30 & 5 (207\,ms) & 3 (450\,ms) & 1 (820\,ms) & 5 (180\,ms) \\
    Detection accuracy (mAP50)   & 0.25 & 4 (0.995)   & 5 (0.997)   & 5 (0.998)   & 3 (0.88) \\
    Model size (TFLite)          & 0.15 & 5 (3.2\,MB) & 3 (22\,MB)  & 2 (50\,MB)  & 5 (4.3\,MB) \\
    TFLite export support        & 0.15 & 5 (native)  & 5 (native)  & 5 (native)  & 4 (manual) \\
    Training data requirement    & 0.15 & 5 (small)   & 4 (moderate)& 3 (large)   & 4 (moderate) \\
    \midrule
    \textbf{Weighted total}      &      & \textbf{4.75} & 3.95       & 3.05        & 4.20 \\
    \bottomrule
  \end{tabularx}
\end{table}

\noindent YOLOv8n is the only variant that achieves real-time inference on the Pi~5 CPU within the 300\,ms-per-frame budget required for useful detection at search speed. Larger variants (YOLOv8s at 11.2\,M parameters, YOLOv8m at 25.9\,M) exceeded this budget by factors of two and four respectively. SSD~MobileNetV2 matches YOLOv8n on speed but delivers lower accuracy on COCO~\cite{lin2014coco}; note that the SSD mAP of 0.88 is measured on COCO whereas the YOLO figures are on our custom dataset, so the comparison is indicative rather than direct~\cite{edgeai}.

The model was trained as a single-class detector on 366 mixed images (300 synthetic, 16 real, 50 negatives) at native camera resolution ($1456 \times 1088$), following the domain-randomisation principle~\cite{domainrand}. Training at higher resolution than the $640 \times 640$ inference input allows the network to learn fine-grained features that survive downscaling~\cite{yolov8}. A COCO person detector is carried as a fallback model for field conditions where the custom model underperforms.

\subsection{Coverage Path Planning}\label{sec:plan-rationale}

Four candidate search patterns were evaluated (Table~\ref{tab:mcda-search}), emphasising coverage guarantee and implementability to ensure reliable, verifiable behaviour over theoretically optimal paths.

\begin{table}[H]
  \centering\compactTable
  \caption{MCDA: search pattern selection (1--5 scale, 5 = best)}
  \label{tab:mcda-search}
  \begin{tabularx}{\textwidth}{X c c c c c}
    \toprule
    \textbf{Criterion} & \textbf{Weight} & \textbf{Lawnmower} & \textbf{Spiral} & \textbf{Expanding Sq.} & \textbf{Random} \\
    \midrule
    Coverage guarantee           & 0.30 & 5 (complete)  & 3 (centre gap) & 4 (approx.)  & 1 (none) \\
    Path length efficiency       & 0.20 & 4 (minimal turns) & 3 (long spiral) & 3 (many turns) & 2 (redundant) \\
    Implementability             & 0.20 & 5 (analytical) & 3 (radius fn)  & 4 (simple)    & 5 (trivial) \\
    Wind robustness              & 0.15 & 4 (long legs)  & 3 (curved)     & 3 (short legs) & 2 (unpredictable) \\
    SAR standard compliance      & 0.15 & 5 (IAMSAR)    & 3 (maritime)   & 4 (land SAR)  & 1 (none) \\
    \midrule
    \textbf{Weighted total}      &      & \textbf{4.65} & 3.00           & 3.65          & 2.15 \\
    \bottomrule
  \end{tabularx}
\end{table}

\noindent The boustrophedon (lawnmower) pattern was selected~\cite{choset2001coverage, galceran2013survey} for its formal coverage guarantee, closed-form path length, and IAMSAR compliance. The implementation uses a rotate--rasterise--rotate-back algorithm on arbitrary polygons, with scan angle auto-optimised via minimum-area bounding rectangle (validated by 216-configuration sweep).

\subsection{State Machine Design}\label{sec:fsm-rationale}

A deterministic finite state machine (FSM) with 20~states and 32~transitions was chosen over behaviour trees or reactive architectures for its verifiability, traceability, and ease of runtime monitoring~\cite{wagner2006fsm}. Each state maps to a single handler function via a dictionary-based dispatch table ($O(1)$ lookup), eliminating fall-through errors. Defence-in-depth safety is enforced through five independent mechanisms: manual override (M~key), RC kill switch (hardware-level), return-to-launch on link loss, three-layer geofence, and per-state timeouts. Every blocking state carries a wall-clock timeout to prevent indefinite hangs. The complete transition table is provided in Appendix~\ref{app:states}.

\subsection{Autonomy Level Justification}\label{sec:autonomy-rationale}

The system operates at Sheridan Level~7--8 (``computer executes autonomously, human monitors'')~\cite{sheridan2002humans} for the search phase and Level~3--4 (``computer suggests one alternative, human approves'') for the landing decision. This hybrid was chosen after analysing the asymmetric cost structure of SAR false positives versus false negatives~\cite{endsley1999loa, cummings2014man_uav}.

A \emph{false positive landing} (descending on a rock, tree stump, or discarded clothing) incurs three compounding penalties: (1)~battery endurance consumed by descent, hover, and re-climb (\SI{\approx 90}{\second}, roughly 8\% of a 20-minute endurance budget); (2)~search coverage paused during the entire descent--verify--climb cycle; (3)~risk of mechanical damage if the landing site is uneven. In a single-drone mission with no redundancy, even one unnecessary landing significantly reduces the probability of covering the full search area before battery exhaustion.

A \emph{false negative} (failing to investigate a real casualty) is worse in absolute terms but is mitigated by the lawnmower pattern's guarantee of revisiting every cell: a missed detection on pass~1 can be recovered on the return leg. The operator-in-the-loop gate therefore addresses the higher-consequence, lower-recoverable error (committing to a landing) while the autonomous search handles the high-bandwidth, lower-consequence task (scanning terrain at 8\,m/s).

The \textsc{Verify} state implements this gate: the operator classifies each detection as confirmed~(Y), rejected~(N), or item of interest~(I) via the ground station. A \SI{120}{\second} timeout auto-rejects unresponded detections---derived from the energy budget (\SI{6.0}{Wh} per hover event, 5.2\% of battery, permitting two timeouts while retaining 89\% for search)~\cite{goodrich2007human}. Full autonomy (Level~10) was rejected because the single-class detector cannot distinguish casualty from visually similar objects~\cite{murphy2014, lyu2023uav_sar_survey}; manual-only operation was rejected because human operators cannot sustain visual scanning at \SI{8}{\metre\per\second} (automation provides ${\sim}20\times$ coverage rate advantage)~\cite{goodrich2008}.

\subsection{Decisions That Changed: Evidence-Based Corrections}\label{sec:decisions-changed}

Four design parameters were revised when empirical evidence contradicted initial assumptions:

\paragraph{Focal length: \SI{7.0}{\milli\metre} $\rightarrow$ \SI{5.46}{\milli\metre}.}
A tape-measure calibration on field day measured \SI{92}{\centi\metre} visible at \SI{1}{\metre} height, revealing a 22\% overestimate in the assumed focal length. This single correction shifted GPS estimations by up to \SI{3}{\metre} at \SI{30}{\metre} altitude---without it, landing errors would have exceeded the \SI{10}{\metre} R07 requirement.

\paragraph{Camera resolution: $640 \times 480$ $\rightarrow$ $1456 \times 1088$.}
At $640 \times 480$, a prone casualty at \SI{30}{\metre} subtends only $20 \times 14$\,px---below reliable detection. Switching to native resolution increased this to $46 \times 32$\,px with no inference time change (the model resizes internally to $640 \times 640$).

\paragraph{TFLite runtime: \texttt{tflite-runtime} $\rightarrow$ \texttt{ai-edge-litert}.}
Google's \texttt{tflite-runtime} has no Python~3.13 wheel; \texttt{ai-edge-litert} provides identical API, eliminating the conflict with \texttt{picamera2} (system Python) with zero code changes.

\paragraph{Serial communication: pyserial $\rightarrow$ MAVProxy UDP bridge.}
A Python~3.13 pyserial regression dropped bytes at 921\,600~baud. MAVProxy resolves this via its C-level serial handler and simultaneously enables multi-consumer telemetry access (Table~\ref{tab:mcda-comms}).

\subsection{Why Not ROS}\label{sec:ros-rationale}

ROS~2 was excluded for three reasons~\cite{quigley2009ros}: (1)~architectural mismatch---the system is a single-node pipeline (one camera, one detector, one flight controller); ROS's pub/sub model adds inter-process overhead without new capability; (2)~deployment complexity---ROS~2 requires $\approx$\SI{2}{\giga\byte} vs four \texttt{pip} packages totalling $<$\SI{50}{\mega\byte}; (3)~timeline---40--60~hr learning curve per developer would have consumed 25\% of available development time for middleware rather than mission logic. ROS remains recommended for future multi-drone extensions.

\subsection{Camera Selection: Global vs Rolling Shutter}\label{sec:camera-rationale}

The IMX296 global shutter camera (provided as university equipment) was validated for aerial search~\cite{imx296, sukhatme2012rolling_global}. At \SI{5}{\metre\per\second} and \SI{30}{\metre} altitude, ground motion across the sensor reaches $\approx$250\,px/s; a rolling-shutter sensor would produce $\approx$7.5\,px parallelogram distortion per frame, shifting detected centroids by up to \SI{0.5}{\metre}. The IMX296's global shutter eliminates this error. The trade-off (lower light sensitivity from \SI{3.45}{\micro\metre} pixel pitch) is acceptable given daylight-only CAA VLOS requirements~\cite{caa2024uas}. During integration, the IMX296 was found to output BGR despite RGB888 labelling---corrected by removing the \texttt{cv2.cvtColor()} call (Section~\ref{sec:decisions-changed}).

\subsection{Field Day Adaptation Under Uncertainty}\label{sec:weather-rationale}

The scheduled flight test (2026-03-11) was cancelled due to sustained winds exceeding the EDU-450's \SI{10}{\metre\per\second} operational limit (measured \SI{12}{\metre\per\second}, gusts to \SI{18}{\metre\per\second}). A pre-planned bench-testing fallback was executed instead, yielding the FOV calibration correction (\SI{7.0}{\milli\metre}~$\rightarrow$~\SI{5.46}{\milli\metre}), lens distortion calibration (RMS~0.399), end-to-end connectivity verification, and the inference benchmark (\SI{207}{\milli\second}, 4.8\,FPS)---results detailed in Section~\ref{sec:decisions-changed} and the evaluation (Section~\ref{sec:evaluation}). The go/no-go decision framework and full bench protocol are documented in Appendix~\ref{sec:field-day}.

\paragraph{Decision flow narrative.} The six-step reasoning chain---from detection ceiling through altitude, speed, scan angle, NFZ margin, and coverage acceptance---is presented as a self-contained executive narrative in Appendix~\ref{app:decision-flow}, including the 216-configuration sweep results, top-three configuration comparison, Pareto optimality analysis, and non-quantifiable design choices.

\medskip
\noindent\textbf{Section summary.} Every design decision in this project traces to a quantified rationale---physical measurement, parametric sweep, weighted trade study, or energy budget analysis---and four decisions were revised when empirical evidence contradicted initial assumptions, demonstrating the evidence-driven engineering process that underpins the system.
